\chapter{Données, protocole et indicateurs}\label{chap:protocole}
\section{Huit nuages de points issus de TSPLIB}
Les fichiers d'entrée proviennent de TSPLIB, bibliothèque de référence pour
le voyageur de commerce\cite{tsplib}. Le projet utilise uniquement les coordonnées
des huit fichiers présents dans le dépôt. Le tableau~\ref{tab:instances} résume
leurs dimensions et leurs types de distance déclarés dans les fichiers sources.
\begin{table}[htbp]\centering
\begin{tabular}{lrrl}
\toprule
Instance & $n$ & $p$ étudiés & Type source \\
\midrule
ulysses16 & 16 & 3, 5, 10 & \code{GEO} \\
ulysses22 & 22 & 3, 5, 10 & \code{GEO} \\
att48 & 48 & 3, 5, 10 & \code{ATT} \\
eil51 & 51 & 3, 5, 10 & \code{EUC_2D} \\
berlin52 & 52 & 3, 5, 10 & \code{EUC_2D} \\
st70 & 70 & 3, 5, 10 & \code{EUC_2D} \\
kroA100 & 100 & 3, 5, 10 & \code{EUC_2D} \\
rd100 & 100 & 3, 5, 10 & \code{EUC_2D} \\
\bottomrule
\end{tabular}

\caption{Instances de la campagne principale. Le type source est conservé comme métadonnée.}
\label{tab:instances}
\sourcecalcul{en-têtes des huit fichiers TSPLIB fournis avec le projet.}
\end{table}

\textbf{La convention de calcul du projet diffère de la métrique officielle
de plusieurs instances.} Toutes les coordonnées sont traitées comme des points
du plan, avec la distance euclidienne non arrondie. La transformation propre
à \code{ATT}, la distance géographique de \code{GEO} et l'arrondi de
\code{EUC_2D} ne sont pas appliqués. Cette convention suit le cadre simplifié
de l'énoncé et rend les méthodes comparables entre elles.

En particulier, les coordonnées des instances Ulysses ne donnent pas ici des
distances physiques en kilomètres. Les optima TSP officiels de TSPLIB ne sont
pas des références valides pour les coûts calculés dans ce projet. Les coûts
bruts de deux instances différentes ne doivent pas non plus être additionnés
ou classés comme s'ils étaient exprimés dans une unité commune.

\section{Plan expérimental réellement exécuté}
Pour chaque instance, trois cardinalités sont retenues : $p=3,5,10$.
Le poids reste fixé à $\alpha=1$, d'où $8\times3=24$ problèmes distincts.
La construction est exécutée une fois. Les recherches locale et tabou sont
répétées sur les graines $0,1,2$, avec 100 itérations chacune. La recherche
tabou utilise une durée de 10 itérations et 15 candidats par itération.

La PLNE n'est tentée que si $n\leq22$ et $p\leq5$, sous une limite de cinq
secondes pour CBC. Elle est exécutée une seule fois dans les cas admissibles,
avec le réglage de solveur enregistré. Les cas qui dépassent les seuils sont
explicitement classés \code{skipped}. Ils ne constituent ni des échecs de
résolution ni des solutions de coût nul.

Au total, le plan contient
\begin{equation}
24\times(1+3+3+1)=192
\end{equation}
enregistrements attendus, en comptant les 20 cas exacts ignorés. Il comporte
172 exécutions avec une solution sauvegardée : 24 constructions, 72 recherches
locales, 72 recherches tabou et quatre PLNE. Parmi ces dernières, trois ont
une preuve et une reste réalisable sans preuve.

Les configurations \code{benchmark_full.json} et
\code{benchmark_time_budget.json} sont proposées pour prolonger l'étude.
La première explore notamment d'autres valeurs de $\alpha$ et davantage de
graines ; la seconde fixe un budget d'amélioration commun. Leur présence
dans le dépôt ne vaut pas résultat expérimental. Le rapport n'attribue donc
aucun effet mesuré au changement de $\alpha$.

\section{Coûts moyens, dispersion et amélioration}
Pour une méthode stochastique et un cas fixé, soient $C_1,\ldots,C_m$ les
coûts des $m$ exécutions réussies. La synthèse conserve notamment la moyenne,
la médiane, le minimum, le maximum et l'écart-type d'échantillon :
\begin{equation}
\overline C=\frac1m\sum_{s=1}^m C_s,\qquad
s_C=\sqrt{\frac{1}{m-1}\sum_{s=1}^m(C_s-\overline C)^2}.
\end{equation}
Avec une seule exécution, aucun écart-type empirique n'est fourni. Avec trois
graines, cet indicateur décrit une dispersion observée limitée ; il n'est pas
une garantie sur l'ensemble des trajectoires possibles et ne constitue pas
un intervalle de confiance.

Le gain face à la construction de coût $C_0$ est
\begin{equation}
G=100\,\frac{C_0-\overline C}{C_0}.
\end{equation}
Il compare la méthode à son point de départ sur le même problème. Il ne mesure
pas la proximité à l'optimum. La synthèse globale des gains donne le même poids
aux 24 cas ; la médiane n'est ni pondérée par $n$ ni calculée sur une somme
de coûts d'unités différentes.

Lorsqu'une valeur optimale $C^*>0$ a été prouvée pour le même cas, on calcule
\begin{equation}
g^*=100\,\frac{\overline C-C^*}{C^*}.
\end{equation}
En l'absence de preuve, cet écart reste non renseigné. La meilleure solution
connue ou une solution PLNE arrêtée au temps limite ne remplace pas $C^*$.

\section{Temps mesuré et environnement}
Le temps d'exécution mesuré inclut le calcul des distances et l'optimisation,
dont l'initialisation des recherches. Il exclut la lecture du fichier source,
l'écriture des résultats et le dessin des figures. Dans le cas PLNE, il inclut
la construction du modèle et son traitement par le solveur. Le temps total
peut donc légèrement dépasser les cinq secondes configurées pour CBC.

Les métadonnées enregistrent Python 3.12.14, PuLP 3.3.2 et Matplotlib 3.11.2,
sous Linux x86\_64. Les versions sont celles de l'exécution sauvegardée, et non
celles d'une documentation en ligne plus récente. Le modèle utilise CBC et
archive son journal. Les durées dépendent de cette machine et du contexte
d'exécution ; les temps très courts de construction sont particulièrement
sensibles au bruit de mesure.

\section{Règles de lecture des figures}
Les comparaisons par instance montrent des coûts et temps moyens, avec les
écarts-types disponibles. La carte des gains compare chaque recherche à la
construction pour le même cas. Les figures de solutions représentent la
meilleure graine sauvegardée et sont identifiées comme telles. Elles ne doivent
pas être interprétées comme une représentation de la solution moyenne.

Les coordonnées conservent leur rapport d'aspect, y compris pour les nuages
allongés. La racine est distinguée, les stations sont visibles et les arcs
d'accès restent plus discrets que l'anneau. L'enjeu est de rendre lisibles
les deux composantes du modèle sans déplacer les points.
\bilan{Les chiffres présentés concernent une campagne précise, avec un poids
unique et trois graines. Les statuts, les unités et l'effort d'évaluation
font partie des résultats au même titre que les coûts.}
